\documentclass[10pt]{article}
\usepackage[margin=20mm]{geometry}
\usepackage{amsmath,amssymb,mathtools,fontspec,unicode-math,xurl,hyperref}
\IfFontExistsTF{DejaVu Serif}{\setmainfont{DejaVu Serif}}{\setmainfont{Latin Modern Roman}}
\setmathfont{latinmodern-math.otf}
\setlength{\parindent}{0pt}\setlength{\parskip}{6pt}\sloppy
\title{Scientific initial review of M25-0020}
\author{}
\date{Review record: 25 October 2026; standard 0.7.1}
\begin{document}
\maketitle
Abstract: This review concerns the 217-page author manuscript arXiv:1710.01722v1 by Mihalis Dafermos and Jonathan Luk, not an established text of the 2025 journal Version of Record. One Major concern is supported in the coordinate-continuity argument of Section 16: the asserted angular-Jacobian limit and transfer of continuity conflict with the manuscript’s own exact rotating Kerr coordinate map. A separate Minor concern concerns degenerate strictly monotone optical initial data in Proposition 2.1; the actual linear initial data of the main theorem avoid it. The first concern invalidates a stated step of the core extension proof in this version, without providing a counterexample to smooth extendibility of Kerr or to the paper’s main physical conclusion. All 217 source pages were viewed, and the supplied TeX and text dependencies were read completely. Independent algebraic checks were executed.\par
\section*{Identity, version and outcome}
Paper. M. Dafermos and J. Luk, The interior of dynamical vacuum black holes I: The \(C^0\)-stability of the Kerr Cauchy horizon. Assigned identifier M25-0020; DOI https://doi.org/10.4007/annals.2025.202.2.1. The bibliographic journal placement is Annals of Mathematics 202 (2025), 309–630. This review is bound to arXiv:1710.01722v1, submitted 4 October 2017, with 5 October 2017 printed on its title page. Both dates are retained; neither is substituted for the other. The 2025 journal text was not compared.\par
Page convention. All source page references below are one-based PDF pages of the acquired 217-page manuscript. Its printed pagination agrees with those page numbers. TeX line references, where supplied, are one-based lines of `KerrInterior.postedversion.tex`.\par
Decision within this review. Problem 1 is Major because a demonstrated invalid coordinate-limit argument is used directly in Theorem 16.14, the final continuity step of the principal extension theorem. Problem 2 is Minor because its degenerate data satisfy the local proposition as stated but are excluded by the main theorem’s actual data. There is no claim of a counterexample to the main nonlinear Einstein theorem, no claim about an unexamined journal version, and no independently accepted scientific decision.\par
Decisive evidence. For exact Kerr with \(M=1\), \(a=3/5\), the manuscript’s angular change is \[\phi_{*,\mathcal{CH}^{+}}=\phi_*-3v,\qquad
 F_{\mathrm K}(v,x)=R(-3v)x\] in either specified stereographic chart. Consequently \(D_xF_{\mathrm K}\) alternates between \(\mathrm I\) and \(-\mathrm I\) along two sequences approaching the same old-coordinate boundary point. The estimate (16.55) would force their difference to tend to zero. This is an exact calculation at admissible rotating subextremal parameters. Problem 2 uses smooth Minkowski annular data \(u=-\tau^3\), \(v=\tau\); the eikonal equation forces \(\mathrm du=0\) at \(\tau=0\), which a coordinate function cannot have.\par
\section*{Claims, hypotheses and coverage}
\subsection*{What the main theorem does and does not assert}
The principal result studied here is the construction, from a spacelike interior hypersurface asymptotic to rotating subextremal Kerr and satisfying the specified weighted smallness assumptions, of a nontrivial Cauchy-horizon portion with a continuous nondegenerate Lorentzian extension. The comparison parameters obey \(0<|a|<M\). This review does not add the Schwarzschild endpoint \(a=0\), the extremal endpoint \(|a|=M\), or uniform constants in either endpoint limit.\par
The introduction also describes event-horizon and bifurcate-horizon formulations and consequences for strong cosmic censorship conditional on further exterior/interior results. Theorems 2 and 3 and the cited subsequent projects are not silently identified with the theorem proved in this manuscript. Likewise, the cosmological-constant discussion on pp. 34–36 is separated from the vacuum \(\Lambda=0\) proof. Section 17 supplies upper regularity estimates; it is not a proof of generic blow-up or failure of unweighted \(H^1\) regularity.\par
The initial parameter is \(\tau\geq1\), with the actual optical data \(v=\tau\) and \(u=C_R-\tau\). The initial weighted norms include metric derivatives through order five and second-fundamental-form derivatives through order four, with polynomial weight \(\tau^{1/2+\delta}\) and \(\delta>0\). The proof uses the detailed initial, null-flux, and sphere norms of Section 4; these are not replaced by a bare informal small-perturbation assumption. Large \(N\) is chosen first, then sufficiently small \(\epsilon\) and sufficiently negative \(u_f\) in the order prescribed in the argument.\par
\subsection*{Basic conventions checked}
Equation reference: review:metric. For the double-null metric \[
 g=-2\Omega^2(\mathrm du\otimes\mathrm dv+\mathrm dv\otimes\mathrm du)
 +\gamma_{AB}(\mathrm d\theta^A-b^A\mathrm dv)
                \otimes(\mathrm d\theta^B-b^B\mathrm dv),\] the manuscript uses \[e_3=\partial_u,\qquad
 e_4=\Omega^{-2}(\partial_v+b^A\partial_{\theta^A}),\qquad
 g(e_3,e_4)=-2.\] A pointwise diagonal angular-metric check gives \(\det g=-4\Omega^4\gamma_{11}\gamma_{22}\) and \(g^{uu}=g^{vv}=0\). The normalization of the unit timelike and spacelike hypersurface normals was checked algebraically, including their squared norms \(-1\), \(1\), and zero inner product. This checks signs and factors, not every component of the Einstein equations.\par
The gauge identities \(\omega=0\), \(\underline\omega=-\partial_u\log\Omega\), \(\zeta=\eta\), \(\nabla\log\Omega=(\eta+\underline\eta)/2\), and \(\partial_u b=2\Omega^2(\eta-\underline\eta)^\sharp\) were followed into their uses. The renormalizations \[K=-\rho+\tfrac12\widehat\chi\mathbin{\cdot}\widehat{\underline\chi}
                  -\tfrac14\operatorname{tr}\chi\operatorname{tr}\underline\chi,
 \qquad
 \check\sigma=\sigma+\tfrac12\widehat{\underline\chi}\wedge\widehat\chi\] are essential to the curvature-pair cancellations; replacing them by unrenormalized components would change the proof.\par
\subsection*{Coverage ledger}
Every source PDF page was viewed. The full 10,977-line main TeX, 660-line bibliography output, 1,097-line bibliography style, and all 24 figure-text overlays were read to their ends. The following ledger describes the mathematical checks, not a claim of independent formal verification of every line.\par
Source pages | Claims and checks\par
1–36 | Introductory theorem scopes, rotating subextremal parameters, conditional consequences, anticipated subsequent work, and cosmological-constant discussion.\par
37–50 | Local optical-coordinate construction, null-frame conventions, null structure/Bianchi systems, and renormalized quantities. Problem 2 is located here.\par
51–65 | Exact initial-data and norm definitions; Theorems 4.2 and 4.24; parameter order; weighted \(L^2\), angular derivative counts, and mixed norm order.\par
65–99 | Bootstrap implications; sphere geometry and analytic estimates; difference equations; highest curvature pairs and derivative cancellation.\par
100–124 | Lower-order estimates and curvature energies; product cases; flux weights; integration by parts and the roles of small parameters.\par
125–140 | Elliptic/transport recovery of third angular derivatives, rescaled incoming quantities, mass aspects, and the coupled \(\underline\omega\) estimate.\par
140–154 | Nonlinear null-norm closure and sphere-norm estimates; near/far splitting; weighted Cauchy–Schwarz and Gronwall.\par
154–165 | Higher regularity at finite upper endpoint; two-dimensional Gronwall; local extension/gluing; construction of the limiting interior region.\par
165–185 | Compactifying coordinate, preliminary continuity, coordinate compatibility, angular Jacobians, tensor-continuity transfer, and the final metric extension. Problem 1 is located here.\par
185–188 | Further weighted regularity estimates and their limits; distinction from unweighted \(H^1\) and generic singularity statements.\par
188–212 | Kerr coordinates and asymptotics; eikonal identity; stereographic charts; explicit rotating horizon map; initial-data and comparison estimates.\par
212–217 | Complete bibliography text and selected identity/version checks; not an independent reading of all cited works.\par
\section*{Problem 1: rotating Kerr contradicts the asserted coordinate limit}
\subsection*{Location and source claim}
Severity: Major. The linked assertions occur in Section 16.1.2, p. 166, especially (16.2); Lemma 16.8, p. 176; (16.44), p. 177; Proposition 16.11, pp. 178–181, including (16.45) and (16.50)–(16.55); Corollary 16.13, p. 182; and Theorem 16.14, pp. 182–183, especially (16.59)–(16.60). The defining Kerr comparison is in Appendix A.5.2, p. 206, with the smooth-extension discussion in A.5.6, p. 207. Related Kerr shift formulas are (A.43), p. 202, and (A.47), p. 203.\par
Equation reference: review:false-tail. The manuscript compares the old compactified coordinates \((u,v_{\mathcal{CH}^{+}},\theta^1,\theta^2)\) with new coordinates \((u,v_{\mathcal{CH}^{+}},\theta^1_{\mathcal{CH}^{+}},\theta^2_{\mathcal{CH}^{+}})\). It asserts the same continuous structure, while allowing a different differentiable structure. Proposition 16.11 asserts both an \(O(\epsilon)\) difference between the actual and Kerr angular Jacobians and continuous extension of the actual angular Jacobian in the old coordinates. Its final estimate (16.55) bounds the same-old-angle difference at two values of \(v\) by a constant times \[
 |v_m^{-2\delta}-v_n^{-2\delta}|^{1/2}
 \left(1+\left\lVert v^{1/2+\delta}\nabla\widetilde b\right\rVert
                _{L^\infty_uL^2_vL^\infty(S)}\right).\] Corollary 16.13 then transfers continuous extendibility of an arbitrary \(S\)-tangent tensor from the old to the new coordinates. Theorem 16.14 explicitly invokes this transfer and the angular-Jacobian limit in proving continuity of the metric.\par
\subsection*{Admissible parameters and all relevant assumptions}
Equation reference: review:kerr-values. Set \[
 M=1,\quad a=\frac35,\quad
 r_-=\frac15,\quad r_+=\frac95,\quad
 \omega_-=\frac{4Mar_-}{(r_-^2+a^2)^2}=3,\quad
 c=\frac{r_+-r_-}{r_-^2+a^2}=4.\] These satisfy \(0<a<M\) strictly. Use exact Kerr initial data on the comparison hypersurface. They solve the vacuum constraints and agree identically with the comparison Kerr data, including all derivatives: the difference size is zero. Therefore every stated smallness requirement on those differences is satisfied for any allowed positive \(\epsilon\), and every positive \(\delta\) is admissible for that comparison. Fix an allowed \(N\), take \(u_f\) as negative as required, and restrict to \(u\leq u_f\). All sequences used below can be started sufficiently late that \(u+v\geq C_R\), the selected initial sphere has been passed, and the exponential formula for \(v_{\mathcal{CH}^{+}}\) applies. Thus neither an early-time boundary nor an excluded parameter endpoint is used.\par
Equation reference: review:kerr-map. In the large-\(v\) region, \(v_{\mathcal{CH}^{+}}=-e^{-4v}/4\), which tends to zero. The manuscript’s Kerr angular map, defined by Appendix A.5.2, is \[
 (\theta_*,\phi_*)\longmapsto
 (\theta_*,\phi_{*,\mathcal{CH}^{+}})=(\theta_*,\phi_*-3v)
 \quad \text{modulo }2\pi.\] An additive initial phase, if retained, merely translates the sequences below and changes none of the conclusions.\par
\subsection*{Exact Jacobian and topological contradiction}
Equation reference: review:rotation, review:sequences. Let \[J_0=\begin{pmatrix}0&-1\\1&0\end{pmatrix},
 \qquad
 R(t)=\begin{pmatrix}\cos t&-\sin t\\ \sin t&\cos t\end{pmatrix}.\] In either of the manuscript’s stereographic charts, the angular coordinate is \[x=\rho(\theta_*)(\cos\phi_*,\sin\phi_*),
 \qquad \rho=\tan(\theta_*/2)\ \text{or}\ \cot(\theta_*/2).\] Consequently the map [review:kerr-map] (\#review:kerr-map) and its derivative are exactly \[
 F_{\mathrm K}(v,x)=R(-3v)x,\qquad D_xF_{\mathrm K}(v,x)=R(-3v).\] Choose \(x_0=(1,0)\) on the equator and the two sequences \[
 v_n=\frac{2\pi n}{3},\qquad
 w_n=\frac{2\pi n+\pi}{3}.\] Then \[F_{\mathrm K}(v_n,x_0)=x_0,\quad F_{\mathrm K}(w_n,x_0)=-x_0,\qquad
 D_xF_{\mathrm K}(v_n,x_0)=\mathrm I,\quad D_xF_{\mathrm K}(w_n,x_0)=-\mathrm I.\] Both images stay in the prescribed sets \(\mathcal V_i'\) with squared radius less than \(2\), and hence in \(\mathcal V_i''\) with squared radius less than \(3\). There is no chart-edge or polar singularity. The Jacobian norm difference is \(2\), and its \((1,1)\) entry has difference \(2\).\par
The old-coordinate sequences \((u,v_{\mathcal{CH}^{+}}(v_n),x_0)\) and \((u,v_{\mathcal{CH}^{+}}(w_n),x_0)\) converge to the same boundary point \((u,0,x_0)\). Their new angular coordinates have distinct limits. Thus the identity transition on the common interior has no continuous extension between these two proposed completions. In particular, they cannot give the same \(C^0\) structure in the sense asserted for that transition. Each completion may be abstractly homeomorphic to a product; that weaker observation does not give the identity-compatible boundary map the proof requires.\par
The right side of [review:false-tail] (\#review:false-tail), for the exact Kerr comparison and any fixed \(\delta>0\), tends to zero, since \[0<v_n^{-2\delta}-w_n^{-2\delta}<v_n^{-2\delta}\longrightarrow0.\] The left side for the Kerr angular Jacobian equals \(2\). This is a direct failure of the supposedly trivial Kerr-background estimate in (16.55).\par
There is a separate ambiguity in the manuscript’s forced coordinate equation, discussed below. The contradiction does not require pretending that its inconsistent stereographic equation has already been solved by [review:kerr-map] (\#review:kerr-map). If its actual solution is used instead, Proposition 16.11 still asserts a uniform componentwise bound \(|D_xF-D_xF_{\mathrm K}|\leq C\epsilon\) and continuity of \(D_xF\). Choose an admissible small \(\epsilon\) with \(C\epsilon<1/2\), using the stated uniform estimate. Along [review:sequences] (\#review:sequences), the \((1,1)\) entries of \(D_xF\) differ by at least \(2-2C\epsilon>1\). They cannot converge to a common boundary value. Exact Kerr satisfies the data hypotheses for this choice; even though \(u_f\) may depend on the smallness choice, the argument applies at that fixed allowed \(u_f\) and arbitrarily large \(v\). The asserted closeness and continuity therefore cannot both hold.\par
\subsection*{Failure of the unrestricted tensor-continuity transfer}
For the intended exact Kerr map, take the globally smooth scalar on the sphere \[f(\theta_*,\phi_*)=\sin\theta_*\cos\phi_*
                 =\frac{2x_1}{1+x_1^2+x_2^2}.\] It is independent of \(u\) and \(v\), so it extends smoothly in the old compactified coordinates. In the new coordinates it is \[f=\sin\theta_*\cos(\phi_{*,\mathcal{CH}^{+}}+3v).\] At \(\theta_*=\pi/2\), \(\phi_{*,\mathcal{CH}^{+}}=0\) its values along [review:sequences] (\#review:sequences) are \(1\) and \(-1\). It has no continuous new-coordinate extension. If the phrase \(S\)-tangent tensor is intended to exclude scalars, use \(f\) times the unit round metric; that metric is rotation-invariant, and the same failure is present in its tensor components. This checks the actual universal continuity-transfer assertion rather than a condition imposed only on axisymmetric metric coefficients.\par
\subsection*{Why the Kerr derivative estimate used in the proof fails}
Equation reference: review:shift-derivative. Write the exact Kerr shift as \(b_{\mathrm K}=B\,\partial_{\phi_*}\), where \[B=\frac{4Mar}{(r^2+a^2)^2-a^2\sin^2\theta\,\Delta},
 \qquad \Delta=r^2-2Mr+a^2.\] At the inner horizon \(B\to\omega_-=3\). In stereographic coordinates this vector field is \(B J_0x\). On compact chart subsets its limiting component derivative is \[
 D_x b_{\mathrm K}\longrightarrow 3J_0\ne0.\] The derivatives of the scalar coefficient \(B\) tend to zero; derivatives of the rotating basis field do not. Formula (A.47) controls special scalar-coordinate derivatives. It cannot be used to conclude that all stereographic component derivatives of \(b_{\mathrm K}\) are \(O(\Omega_{\mathrm K}^2)\).\par
The estimate used before (16.51), p. 181, consequently does not give \(\int^\infty |D_xb_{\mathrm K}|\,\mathrm dv<\infty\): already an off-diagonal component tends to \(3\) in absolute value. This also explains why treating the background rotation as a negligible transport coefficient is invalid. A rotation can preserve geometric distances without having a small componentwise Jacobian generator; the two statements are different.\par
\subsection*{Chart consistency requires the derivative at the image point}
There is a related coordinate issue in Lemma 16.8 and equations (16.39)–(16.44). If an angular map \(F\) satisfies \[(\partial_v+b\mathbin{\cdot}\partial_x)F=q(x),\] then a nonlinear target-chart change \(\Psi\) gives \[(\partial_v+b\mathbin{\cdot}\partial_x)(\Psi\circ F)
       =D\Psi(F(x))q(x).\] The components of the old-point vector field in the changed chart instead use \(D\Psi(x)q(x)\). The proof equates these derivative evaluations without justification. The elementary local diagnostic \(F(v,x)=x+v\), \(\Psi(x)=x^2\) on \(x>0\) gives \(2(x+v)\) on the left and \(2x\) for the asserted transformed forcing. This diagnostic is only a chain-rule check; the Kerr calculation supplies the counterevidence in the relevant spacetime.\par
Equation reference: review:covariance-residual. Even granting the natural correction of the subtracted constant in (16.40) to \(\omega_-\), the polar Kerr equation reads \[(\partial_v+B\partial_{\phi_*})\phi_{*,\mathcal{CH}^{+}}=B-\omega_-.\] For the stereographic Kerr map [review:rotation] (\#review:rotation), its correct left side is \((B-\omega_-)J_0R(-\omega_-v)x\). The stated old-point forcing in (16.44) is \((B-\omega_-)J_0x\). Their difference is \[
 (B-\omega_-)(R(-\omega_-v)-\mathrm I)J_0x,\] which is generally nonzero before the horizon. At the equator with the parameters [review:kerr-values] (\#review:kerr-values), simplification gives \[B(r)=\frac{12/5}{r^3+(9/25)r+18/25},\qquad r_-<r<r_+.\] Its denominator is strictly increasing for \(r>0\) and gives \(B(r_-) =3\); hence \(B(r)<3\) for every interior \(r>r_-\). For any fixed allowed \(u_f\), late finite \(v\) has \(r>r_-\), and infinitely many such \(v\) have a nonidentity rotation. The residual is therefore nonzero at legitimate late points. For illustration, \(r=1/2\) gives \(B=96/41\), and algebraic substitution of \(R=-\mathrm I\), \(x=(1,0)\) gives residual \((0,54/41)\). This last numerical pair is a finite algebra illustration, not an additional assertion that a prescribed fixed-\(u_f\) orbit simultaneously takes those illustrative values.\par
\subsection*{Counterevidence and precise impact}
Several facts limit the conclusion. First, the exact Kerr spacetime is smoothly extendible in the regular rotating coordinates described in the appendix; the calculation above does not refute that. Second, \(R(-3v)\) is invertible with determinant \(1\) for every finite \(v\). Failure of its boundary limit is not a failure of interior diffeomorphism. Third, the Kerr metric has axisymmetry, which can cancel rotations in transformed metric components even when a generic tensor does not extend. Fourth, local polar angular derivatives can hide the rotating phase, but they do not remove it or supply a globally regular stereographic boundary map. Reversing the rotation sign or using the inverse map leaves the nonconvergence unchanged.\par
The dependency is nevertheless central: \[\begin{gathered}
 \text{Proposition 16.11}
 \ \Longrightarrow\ \text{Corollary 16.13}
 \ \Longrightarrow\ \text{Theorem 16.14}\\
 \Longrightarrow\ \text{the stated continuous extension}.
\end{gathered}\] Theorem 16.14 specifically transfers continuity and uses inverse angular-Jacobian continuity. The source’s proposed proof of those prerequisites contains the exact contradiction above. A replacement argument is required for the nonlinear perturbed metric in the reviewed version. The significance does not depend on showing that replacement is impossible.\par
Proposition 16.6’s conversion of transport control into fixed-old-angle continuity also deserves care because the relevant vector contains an angular shift. This report does not count that observation as a further demonstrated counterexample to its metric conclusion: exact Kerr’s symmetry gives genuine cancellations. Likewise, the ordinary matrix exponential written in the proof of Proposition 16.11 would require commutation or a correctly ordered fundamental solution; it is not counted as an additional independent issue here.\par
\subsection*{Supported repair ingredient and remaining nonlinear work}
A precise limited replacement ingredient is available, but it is not a completed repair of Section 16.\par
Lemma 1.  Let \(V\) generate a fixed round-sphere rotation, let \(\omega\) be constant, and let \(b(v)=\omega V+q(v)\) be a smooth time-dependent vector field on \(\mathbb S^2\). Suppose \(\int_{v_0}^\infty \left\lVert q(v)\right\rVert_{C^1(\mathbb S^2)}\,\mathrm dv<\infty\). If \(\Phi_v\) is its forward flow from \(v_0\) and \(Q_t\) is the rotation flow of \(V\), then \(Q_{-\omega(v-v_0)}\circ\Phi_v\) converges in \(C^1\) to a \(C^1\) diffeomorphism.\par
Proof. Conjugating by the round rotation produces a flow \(\Psi_v\) with vector field \[\widetilde q_v=(Q_{-\omega(v-v_0)})_*q_v.\] Rotations preserve the relevant norms up to a fixed equivalent choice of \(C^1\) norm. Thus \[\int\left\lVert\widetilde q_v\right\rVert_{C^1}\,\mathrm dv<\infty.\] Its displacement between times \(v_1\) and \(v_2\) is bounded by \(\int_{v_1}^{v_2}\left\lVert\widetilde q_v\right\rVert_{C^0}\,\mathrm dv\). The derivative equation is \(\partial_vD\Psi_v=D\widetilde q_v(\Psi_v)D\Psi_v\). Gronwall gives a uniform derivative bound; integration of that equation gives a uniform Cauchy estimate for \(D\Psi_v\) with the \(C^1\) tail integral. The inverse derivatives are uniformly bounded by the same integrable-coefficient argument. Hence the limit is locally invertible and globally bi-Lipschitz. Surjectivity follows by compactness: for each target, take a convergent subsequence of its preimages under \(\Psi_v\). The limit and inverse are \(C^1\) by local inversion, giving the assertion. ◻\par
This lemma removes the persistent rigid rotation before asking for convergence. To use it in the full problem one must construct an intrinsic compatible sphere map, establish the needed bounds in the precise mixed norms and on the relevant hypersurfaces, and prove continuity of the transformed metric directly or through a valid relative-coordinate transfer. Lemma 16.8, the applicable part of Proposition 16.6, Proposition 16.11, Corollary 16.13, and Theorem 16.14 would need corresponding revisions. Those nonlinear revisions have not been proved in this review. Simply deleting the nondecaying rotation from one displayed estimate does not verify the original conclusion.\par
\section*{Problem 2: strict monotonicity does not prevent a degenerate optical coordinate}
\subsection*{Location, claim and severity}
Severity: Minor. Proposition 2.1, pp. 38–39, constructs local double-null coordinates near the embedded spacelike hypersurface \(\Sigma=\ell\times\mathbb S^2\). The simple curve \(\ell:(a,b)\to\mathbb R^2\) supplies smooth functions \(u|_\Sigma\) and \(v|_\Sigma\) that are assumed strictly decreasing and strictly increasing, respectively, in its parameter \(\tau\). The proposed coordinates satisfy the eikonal equations and metric form (2.1). The stated strict monotonicity, in its ordinary mathematical meaning, does not imply nonzero derivatives.\par
\subsection*{Complete counterexample to the stated local proposition}
Equation reference: review:degenerate-curve. Consider the annular region of Minkowski spacetime \[g=-\mathrm dt^2+\mathrm d\tau^2+(2+\tau)^2g_{\mathbb S^2},
 \qquad -1<\tau<1,\] where \(g_{\mathbb S^2}\) is the unit round metric and the ordinary spatial radius is \(2+\tau>0\). Take \(\Sigma=\{t=0\}\) with all of \(\mathbb S^2\), and \[
 \ell(\tau)=(-\tau^3,\tau),\qquad
 u|_\Sigma=-\tau^3,\qquad v|_\Sigma=\tau.\] The ambient metric is smooth and Lorentzian. The induced metric on \(\Sigma\) is positive definite, so the embedding is spacelike and isometric with the specified induced metric. The curve is smooth and injective; its derivative \((-3\tau^2,1)\) never vanishes, so it is a regular simple curve with smooth inverse parameter \(v\). The function \(-\tau^3\) is strictly decreasing and \(\tau\) strictly increasing on the whole interval. Thus this is not a nonsmooth parameterization or a curve cusp. Additive constants may be added to \(u\) and \(v\) if a coordinate-range translation is desired.\par
At any point of the sphere \(\tau=0\), all tangential derivatives of \(u|_\Sigma\) vanish. With \(N=\partial_t\) the unit normal, any smooth extension of \(u\) has \[g^{-1}(\mathrm du,\mathrm du)
       =-(Nu)^2+|\mathrm d_\Sigma u|^2
       =-(Nu)^2\] at that point. The required eikonal equation therefore gives \(Nu=0\). Hence \(\mathrm du=0\) as a full spacetime covector at that point. A function whose differential vanishes cannot be one member of a coordinate system there. In particular, it cannot give the asserted nonsingular double-null metric with positive finite lapse.\par
The proposition asserts the existence of such coordinates near every point (and the associated sphere neighborhood); the point \(\tau=0\) contradicts that assertion. The counterexample satisfies every stated geometric and monotonicity hypothesis, while the asserted conclusion is impossible.\par
\subsection*{Alternative reading, supported correction and limited impact}
If strictly decreasing/increasing was intended to mean that the derivatives are strictly negative/positive everywhere, that is the missing explicit assumption and excludes [review:degenerate-curve] (\#review:degenerate-curve). These meanings are not equivalent for smooth real functions.\par
The precise local correction is to require \[\partial_\tau u|_\Sigma<0,\qquad
 \partial_\tau v|_\Sigma>0\] pointwise, in addition to the existing smooth spacelike hypotheses. The functions are constant on each initial sphere. If \(n\) is the unit normal to those spheres within \(\Sigma\), oriented by \(n\tau>0\), then \(nu<0\) and \(nv>0\). Choosing the appropriate eikonal branches gives nonzero \(Nu\) and \(Nv\). The normal derivative of the eikonal equation, viewed as an equation for the normal first derivative, is \(-2Nu\) or \(-2Nv\) and is nonzero. Thus the ordinary smooth noncharacteristic local construction applies, with the prescribed branch; compactness of the sphere permits a common sufficiently small neighborhood. This verifies the missing nondegeneracy ingredient of the proposition, not a new global existence statement.\par
The actual main data (4.5), p. 52, are \(u=C_R-\tau\) and \(v=\tau\). Their derivatives are \(-1\) and \(1\). Equations (4.8)–(4.10), p. 55, have the positive factor \(n\tau=(\Phi-|w|^2)^{-1/2}\) under the spacelike condition. Accordingly the degenerate example does not invalidate the invocation for these main data. This supported defect remains local and Minor. The reciprocal normalization appearing in the nearby proof is not counted as another issue: the geometric normal calculation and later formulas make the intended normalization recoverable, and isolated typographical defects are outside this review’s issue count.\par
\section*{Other checked reasoning and scope resolutions}
\subsection*{Weights, transport and derivative count}
The comparison asymptotic is \[\Omega_{\mathrm K}^2\asymp e^{-2\kappa_-(u+v)},\qquad
 2\kappa_-=\frac{r_+-r_-}{r_-^2+a^2}.\] The incoming quantities \(\underline\chi_{\mathrm K}\), \(\underline\omega_{\mathrm K}-\kappa_-\) and \(\underline\beta_{\mathrm K}\) carry the small \(\Omega_{\mathrm K}^2\) factor, whereas the corresponding outgoing or angular quantities need not. The review retained these distinctions when following the weighted initial and flux norms. In particular, the order of the \(L^\infty\) and \(L^2\) norms at the top angular order is material and cannot be freely interchanged.\par
Equation reference: review:tail. For \(\delta>0\) the elementary tail estimate \[
 \int_L^\infty v^{-1-2\delta}\,\mathrm dv
      =\frac{L^{-2\delta}}{2\delta}\] justifies an \(L^1\) transport bound from a weighted \(L^2\) flux by Cauchy–Schwarz. At \(\delta=0\) that integral diverges. The proof’s positive \(\delta\) assumption was retained throughout; no endpoint extension is asserted here. This valid tail estimate does not control a nondecaying background rotation as though it were a small difference field.\par
The auxiliary weight \(\varpi=1+e^{-2\kappa_-(u+v-C_R)}\) lies in \([1,2]\) in the interior domain and has derivative equal to \(-2\kappa_-\) times its exponential part. Its favorable differentiated term and the comparison between \(\Omega\) and \(\Omega_{\mathrm K}\) were followed in the energy estimates.\par
\subsection*{Highest curvature equations and elliptic recovery}
The paired equations of Section 7.7, pp. 96–99, are organized so that the top angular derivatives of \(K\) and \(\check\sigma\) cancel with the divergence terms of the differentiated \(\beta\) or \(\underline\beta\) equations after integration by parts. For example, the outgoing pair has schematic principal part \[\nabla_3\nabla^2(\Omega\Omega_{\mathrm K}\widetilde\beta)
    +\nabla(\Omega\Omega_{\mathrm K}\nabla^2\widetilde K)
    -{}^*\nabla(\Omega\Omega_{\mathrm K}\nabla^2\widetilde{\check\sigma})
        =\text{controlled errors},\] together with \(\nabla_4\nabla^2\widetilde K+
\operatorname{div}\nabla^2\widetilde\beta\) and the dual-divergence equation for \(\widetilde{\check\sigma}\). The factor \(\Omega\Omega_{\mathrm K}\) matches the associated \(\Omega^2\Omega_{\mathrm K}^2\) flux and is not an arbitrary decoration.\par
For angular products through total order three, the checked estimates separate the zero-derivative \(L^\infty\) case, the low-order \(L^\infty\) versus \(L^2\) case, and the \(L^4\)–\(L^4\) distribution of derivatives. This avoids placing an uncontrolled third derivative in \(L^\infty\). The elliptic/transport recovery of third Ricci derivatives, the rescaling of \(\underline\chi\) by \(\Omega_{\mathrm K}^{-2}\), and the mass-aspect equations were traced into Section 11. The \(\underline\omega^\dagger\) combination and the incoming curvature equation provide the coupled control used for the third \(\underline\omega\) derivative.\par
The check of Proposition 11.5’s product bound retained the available \(L^2\) control of the metric difference when the other factor is controlled by the lower-order Sobolev estimates. It did not substitute an independent metric bound not supplied by the norms. The coupled estimates in Proposition 11.8 use the small \(N^{-1}\) factors in their designated order.\par
These are checks of the structure, weights, derivative distributions, and cited closure paths. They are not a line-by-line independent rederivation of every commutator coefficient in the 217-page argument. No additional supported defect is asserted from a coefficient whose full alternative derivation was not completed.\par
\subsection*{Null-norm closure and finite-endpoint continuation}
Section 12 splits the spacetime error region into a portion where \(|u|\) is comparable to \(v\) and a portion with stronger exponential decay. The first uses the polynomial weights; the second uses the \(\Omega\) factors after taking \(u_f\) sufficiently negative. The resulting \(N^{-1}\) factor and the higher-power small-norm errors are compatible with the stated bootstrap order. The closure uses continuity on the small branch, not merely an algebraic inequality with a large admissible branch.\par
The abbreviated \(\mathcal N_{H,1}\) bookkeeping in Section 12 does not list every quantity appearing in all subsequent errors, notably \(\underline\omega\). The already available \(\underline\omega\) flux and elliptic estimates, including Proposition 9.11 and Proposition 11.8, allow it to be included in the controlled auxiliary collection. This was recorded as a bookkeeping clarification rather than a separate demonstrated core failure.\par
For the higher-regularity continuation, constants are allowed to depend on a finite upper endpoint \(v_f\). They are not silently claimed uniform at the Cauchy horizon. The two-dimensional Gronwall calculation of Lemma 14.1 was checked by integrating \[\int_0^q 2A e^{8Bt}\,\mathrm dt
       =\frac{A}{4B}(e^{8Bq}-1).\] The two feedback integrals multiplied by \(B\) contribute at most \(A e^{8Bq}/2\), which gives the strict bootstrap improvement \(A+(A/2)e^{8Bq}\leq (3A/2)e^{8Bq}<2Ae^{8Bq}\) for \(q\geq0\). The angular and spacetime derivative inductions and the use of smooth local Cauchy/characteristic uniqueness in the nested union were followed. The underlying general local-existence theorems were used as cited dependencies, not independently re-proved.\par
\subsection*{Mechanical discrepancies not promoted to substantive issues}
Several apparent discrepancies were resolved locally or classified as isolated notation errors, rather than counted as supported problems:\par
- The tail of the initial hypersurface corresponding to \(u<u_f\) has \(\tau>C_R-u_f\). The reversed inequality in some displayed descriptions is contradicted by the coordinate definition and the surrounding tail interpretation. It is not used here to manufacture an empty-domain counterexample.\par
- In the cubic energy estimates around pp. 115–116, an isolated displayed power of \(\Omega\) must be read with the adjacent weight decomposition and the later correctly weighted estimate (10.43), p. 122. The actual product estimate can retain the required cubic factor.\par
- The one-sided squared-norm comparison around (10.28) must include its cross term. Young’s inequality bounds it using the small energy parameter and the already controlled lower curvature term; it does not by itself establish a new core failure.\par
- The \(b\) sphere estimate has a rough square-root bound in an intermediate presentation, while its linear transport equation with integrable coefficients gives the required small bound from the small initial value and source. The stronger bound is the one needed for closure.\par
- The conformal-lapse renormalization in Section 16.2 uses the \(\Omega^2\) transport coefficient. An isolated \(\Omega\Omega_{\mathrm K}\) occurrence was not treated as an independent issue when the actual renormalized equation and equivalent weighted flux supply the intended factor.\par
Each resolution is limited to the displayed bookkeeping issue. It does not repair Problem 1.\par
\subsection*{Kerr asymptotics and the regularity statement}
The Appendix A eikonal algebra was checked by substituting the definitions of \(\Sigma\), \(R\), and \(\Delta\). The relevant rational-expression residual is exactly zero. For the near-horizon exponential asymptotic, mere comparability of logarithms would not alone establish the exact exponential rate. The underlying radial equation provides the stronger argument: if \(y=r-r_->0\) satisfies \(y'=-cy+O(y^2)\), a preliminary exponential decay estimate makes \(\int y\) finite; then \((\log y+cr_*)'=O(y)\) is integrable, so \(y\asymp e^{-cr_*}\). This supplies the needed inference, with the actual radial orientation.\par
The negative-\(a\) orientation in the angular integral formula is handled by the Kerr symmetry \(a\mapsto-a\), \(\phi\mapsto-\phi\) if one first fixes \(a>0\) for that construction. The review does not count that removable orientation convention as a separate Major or Minor claim, and the counterexample in Problem 1 uses \(a>0\) directly.\par
For Section 17, put \(s=-v_{\mathcal{CH}^{+}}>0\) near the horizon. At fixed finite \(u\), \[\Omega_{\mathrm K}^2\asymp s,\qquad
 |\Omega_{\mathrm K}^4\,\mathrm dv|\asymp s\,\mathrm ds.\] Thus a bound involving \(s(\log(1/s))^{1+2\delta}|\chi|^2\) is a weighted integrability assertion. It does not imply the corresponding unweighted \(H^1\) bound, and an upper bound of this kind does not prove generic blow-up. The signs and absolute values in the logarithmic coordinate are interpreted using this positive variable.\par
\section*{Dependence, version limits and conclusion}
Problem 1 requires a replacement for a core continuity argument in the bound 2017 manuscript. No fresh empirical experiment is needed to verify the exact rotation contradiction; symbolic algebra and the explicit map suffice. Validating a repaired nonlinear extension proof would require new analytic work, not an Einstein evolution simulation. Problem 2 is corrected at the level of local noncharacteristic assumptions and does not reach the actual linear main-data construction.\par
The primary arXiv identity page and publisher metadata were checked. Bounded searches for a correction matching the paper and Proposition 16.11 did not establish a matching erratum or a revised proof. This is a search limitation, not evidence that no correction exists. In particular, the longer 2025 journal article was not matched section by section to this manuscript. No journal-final defect is claimed.\par
Item | Judgment | Repair, adoption and remaining work\par
Problem 1 | Major; core derivation invalid in v1 | Exact rotating-Kerr check complete. Relative-flow lemma proved; full nonlinear repair unverified. Author adoption not requested; independent acceptance pending.\par
Problem 2 | Minor; local hypothesis insufficient | Counterexample complete. Pointwise derivative condition restores local nondegeneracy; main linear data already satisfy it. Adoption not requested.\par
Other checks | No additional supported issue within the stated checks | Scope clarifications retained; no blanket formal-correctness claim.\par
\section*{Supporting checks}
The reviewer-owned Python program\par
> `[local archive path omitted]\par
with the `-X utf8 -B` options. It did not import or run an author program. The successful recorded process was PID 51788, exit code 0, with a recorded runtime of approximately 3.02 seconds. The result timestamp is `2026-10-25T23:50:18.843977+00:00`. The runtime reported Python 3.13.5 and SymPy 1.13.3.\par
Thirteen check groups returned true:\par
1. Exact rotating subextremal Kerr parameters and zero comparison-data difference.\par
2. The two rotating sequences, chart containment, map values and Jacobian gap.\par
3. The vanishing bound in (16.55), by an analytic positive-tail inequality, with a finite algebra check at \(\delta=1/2\).\par
4. The old/new scalar-continuity example and tensor alternative.\par
5. The nonzero limiting stereographic derivative of the Kerr shift.\par
6. The Kerr componentwise coordinate-equation residual, including the stated illustrative rational value.\par
7. The nonlinear chart derivative evaluation-point diagnostic.\par
8. The regular strictly monotone initial curve and vanishing optical covector.\par
9. Null-frame, determinant and unit-normal algebra.\par
10. The two-dimensional Gronwall integral.\par
11. The weighted tail primitive and the positive-\(\delta\) condition.\par
12. The Kruskal variable and singular-weight measure conversion.\par
13. The Appendix A eikonal rational identity.\par
The analytic arguments in the report state the limiting reasons; a Boolean result is not used as a substitute for them. These are exact algebra and elementary analytic checks, not a numerical nonlinear evolution, formal proof certification, or experimental reproduction.\par
Failed attempt retained. An earlier reviewer-script version stalled in a generic symbolic limit calculation. Its PowerShell process PID 11380 and Python child PID 28824 were identified by their actual command line; only that reviewer-owned Python child was stopped. The attempted script is retained as `reviewer\_checks\_attempt1.py`. The successful revision replaced the problematic generic symbolic limit by the displayed analytic inequality and a finite algebra check, and used a differentiated primitive for the tail integral. No failed result is represented as a passing computation.\par
\section*{Source inspection, hashes and rendering evidence}
The input manifest was read completely (768 lines; zero remaining), and all 58 files bound by it were rehashed with no mismatch in the recorded batch integrity check. The main evidence identities are:\par
- Acquired PDF, 2,899,501 bytes; SHA-256
`e3d35d61cee99bc0dcd6f0574ccce09497be`
`89f163feb5a0584791da15591eef`.\par
- Main TeX, 893,636 bytes; SHA-256
`56acd2fb4a53fc7f0a1b414fbb303fb4c`
`870d38184f979ec5bdf02f83f2ed0b5`.\par
- Manifest SHA-256
`b6b6287a5e0253bcefb11bc79a7e70aff`
`06e8eec46dc1cece64a6dda523ebf6c`.\par
- Successful check-script SHA-256
`befafbde98f0e8e81eb11bea75e3043f`
`66099ca2c958ae506c6816deea7011e5`.\par
Line wrapping above does not insert characters into the hashes; the exact uninterrupted values and every input-member hash are in the JSON provenance and audit records.\par
Poppler `pdftotext` extraction completed in PID 41960, exit code 0, and produced all 217 page entries. Poppler `pdftoppm` rendering completed in PID 41436, exit code 0, producing `source\_pages/page-001.png` through `page-217.png`. Every one of those page images was actually opened and viewed by this assistant. The renderer emitted `Syntax Error: No display font for ’Symbol’`. The resulting page images remained legible; equations used for the findings were also checked against the complete source text. The warning is preserved rather than hidden.\par
Static source inspection inventoried 51 supplied source-tree members. The main file has 784 labels and 758 reference occurrences, and the supplied bibliography has 115 keys corresponding to 115 citation keys in the audit; the regex check found no unresolved reference or citation candidates. The scanned main-source patterns for `\textbackslash{}write18`, `\textbackslash{}directlua` and `\textbackslash{}openout` were absent. This is a bounded static scan, not a TeX interpreter or a security certification. The named original `.bib` was not bundled, but the complete `.bbl` was available and read. No numerical author program was present in the inspected bundle.\par
Useful one-based TeX anchors are 9069 (differential structures), 9319 (preliminary continuity), 9451 (forced coordinate equation), 9467 (coordinate consistency), 9508 (Kerr transport equation), 9560 (angular-Jacobian proposition), 9680 (the estimate (16.55)), 9726 (tensor transfer), and 9734 (metric extension). The full source, exact page text and rendered images permit independent recovery of each location.\par
\section*{Supporting checks}
Only its assigned Windows staging directory was used for output. The evidence directory contains extraction/page records, all source-page images, the reviewer scripts and results, static audit, inspection receipt, external-reference check record, compilation status and delivery inventory.\par
The compiler capability record found no `pdflatex`, `xelatex`, `lualatex` or `tectonic` on the checked PATH and standard locations. No compiler was installed. Neither the report nor the author TeX was compiled. Accordingly `reports/Pdf/0020.pdf` does not exist as a produced report, and no report-page visual inspection is claimed. A later environment-balance check is only static lint; it is not compilation. The recoverable scientific report is saved, but the full delivery flag remains false.\par
All scientific reading, mathematical reasoning, reviewer computation, static inspection and actual source-page viewing in this review were performed by this assistant. No scientific subagent was used. Human scientific approval, author adoption and independent acceptance are not claimed.\par
This observation belongs to the dispatcher. This assistant did not personally observe or take control of the dispatcher-held tab. No equivalence between that runtime name and the UI label is asserted; the backend route was not independently verified. No independent evidence establishing a non-6-Pro route was observed, and no model or mode switch was performed.\par
General local-existence theorems, all cited historical results and every lower-order commutator coefficient were not independently re-proved. The claim-focused reading and explicit checks support the stated findings with the limitations recorded above; they do not certify the entire manuscript. The 2025 journal Version of Record was not compared. A separate substantive acceptance review and the missing genuine compilation/page checks remain necessary.\par
9 M. Dafermos and J. Luk, The interior of dynamical vacuum black holes I: The \(C^0\)-stability of the Kerr Cauchy horizon, arXiv:1710.01722v1, submitted 4 October 2017, title-page date 5 October 2017, 217 pages. Exact acquired manuscript reviewed. https://arxiv.org/abs/1710.01722v1.\par
arXiv author-repository identity and version record for 1710.01722, consulted 25 October 2026. https://arxiv.org/abs/1710.01722. Used for identity/version metadata, not as proof of any unexamined revision.\par
Annals of Mathematics, publisher identity record for Dafermos–Luk, vol. 202 (2025), no. 2, pp. 309–630, DOI 10.4007/annals.2025.202.2.1, consulted 25 October 2026. https://annals.math.princeton.edu/2025/202-2/p01. Journal-final text correspondence not established.\par
\section*{Compilation continuation (26 October 2026)}
Scientific conclusions, pending questions, source-version limits and the absence of independent acceptance remain unchanged.\par
\end{document}
